\documentclass[../../main.tex]{subfiles}

\begin{document}

\section{Il progetto di DUNE} \label{sec:il-progetto-di-dune}
  \begin{figure}[ht]
    \includegraphics[width=\textwidth]{dune}
    \caption{Illustrazione del progetto per l'esperimento DUNE}
    \label{fig:illustrazione-dune}
  \end{figure}
  DUNE sarà un esperimento a lunga distanza che mira a osservare l'oscillazione
  dei neutrini.
  Come illustrato nella figura~\ref{fig:illustrazione-dune}, l'esperimento sarà
  composto da una nuova sorgente di neutrini ad alta intensità, prodotta da
  un acceleratore di protoni con potenza dell'ordine del megawatt operante al
  Fermilab, un detector composito collocato vicino alla sorgete, d'ora in
  avanti ``\textit{near detector}'' (ND), e un secondo detector molto più grande
  del primo posizionato 1,5 km sotto terra a 1300 km di distanza dal ND presso
  lo Sanford Underground Research Facility, d'ora in avanti
  ``\textit{far detector}'' (FD).

\end{document}